\section{Hamiltonian comparison on a toric domain}\label{cmp:section}

Our goal is to turn a positive mean area surplus into a strict
pointwise inequality after a Hamiltonian change of fiber coordinates.
The construction must also preserve the boundary moments needed by
the surface deformation in Section~\ref{surf:section}.

Let
\begin{equation}\label{cmp:polytope}
 \Delta=\{p\in\R_{\ge0}^m:
             b_\nu\cdot p\le c_\nu,\ 1\le\nu\le N\},
 \qquad b_\nu\in\R_{\ge0}^m,\quad c_\nu>0,
\end{equation}
be compact and full dimensional, with \(m\ge1\).
We call the nonempty faces \(b_\nu\cdot p=c_\nu\) the
\emph{outer faces}. The coordinate faces \(p_j=0\) are not outer faces.
Put
\[
 V=\{z\in\C^m:p(z)\in\Delta\},\qquad
 p_j(z)=\pi|z_j|^2,
\]
with its standard symplectic form \(\omega_V\).
We freely regard functions on \(\Delta\) as toric functions on \(V\).
Smooth maps, Hamiltonians, and differential forms on these compact
sets will always be defined on neighborhoods. No simplicity or
rationality assumption on \(\Delta\) is needed.

\begin{lemma}[Hamiltonian comparison]\label{cmp:comparison}
Let \(P:\Delta\to\R\) be continuous and concave. Suppose that its
Lebesgue mean \(\overline P\) is positive and, for some \(0<\tau<1\),
\begin{equation}\label{cmp:outer-positive}
 \inf_{\Delta\setminus\tau\Delta}P>0.
\end{equation}
Then there are a smooth toric function \(K\) and a Hamiltonian
diffeomorphism \(q:V\to V\) such that
\begin{equation}\label{cmp:conclusion}
 (K-P)(qu)<K(u)\qquad(u\in V).
\end{equation}
The map \(q\) admits a Hamiltonian isotopy from the identity preserving
\(V\), whose Hamiltonians Poisson commute with \(b_\nu\cdot p\) near
each corresponding outer face, uniformly in time.
\end{lemma}

We separate the rearrangement inequality from its symplectic realization.
First we define the rearrangement operator and state the realization
property needed by the proof.  The fixed-point argument then produces a
uniform gap; the remaining subsection proves the realization property
with an arbitrarily small error and the required boundary control.

\subsection{Coordinate rearrangements}

For a continuous convex function \(f:[0,1]\to\R\), its increasing
rearrangement with respect to Lebesgue measure at \(h\in[0,1]\)
can be written as
\begin{equation}\label{cmp:interval-formula}
 g(h)=\min_{0\le a\le1-h}\max\{f(a),f(a+h)\}.
\end{equation}
This also defines the endpoint values \(g(0)=\min f\) and
\(g(1)=\max f\).

For a continuous convex function \(F\) on \(\Delta\), let \(R_jF\)
denote increasing rearrangement in \(p_j\), with the other moments
\(y=(p_i)_{i\ne j}\) fixed. The slice is \([0,L(y)]\), where
\[
 L(y)=
 \min_{\nu:b_{\nu j}>0}
 \frac{c_\nu-\sum_{i\ne j}b_{\nu i}y_i}{b_{\nu j}}.
\]
At least one such bound exists by compactness of \(\Delta\).
The function \(L\) is continuous, concave, and piecewise affine
on the projected polytope. Define \(R_jF\) by the rescaled version
of~\eqref{cmp:interval-formula}.

\begin{lemma}\label{cmp:rearrangement-properties}
Each \(R_j\) maps continuous convex functions on \(\Delta\) to
continuous convex functions. It preserves Lebesgue integrals,
order, and addition of constants, and is nonexpansive in uniform
norm. Its output is nondecreasing in \(p_j\); moreover it preserves
being nondecreasing in any other coordinate.
Consequently \(R=R_m\cdots R_1\) has coordinatewise nondecreasing
output.
\end{lemma}

\begin{proof}
Continuity follows from the minimum-over-intervals formula. At a
slice of length tending to zero, uniform continuity of \(F\) gives
the same conclusion. For convexity, choose minimizing intervals
of lengths \(h_0,h_1\) on two slices. Their convex combination is
an allowed interval of length \((1-t)h_0+th_1\) on the intermediate
slice. Convexity of \(F\) bounds its two endpoint values by the
corresponding convex combinations, which proves convexity of
the rearrangement.

The minimum formula is nondecreasing in the interval length.
If \(F\) is nondecreasing in another coordinate, lowering that
coordinate preserves every allowed interval, by the nonnegative
coefficients in~\eqref{cmp:polytope}, and can only lower its
endpoint values. This proves preservation of the previous
monotonicity properties.

On each slice, increasing rearrangement preserves the distribution
of values, including flat sublevels, and hence the integral.
Integrate this identity in the other coordinates.
Order and translation invariance follow directly from the
minimum formula. Applying these to
\(G-\norm{F-G}_\infty\le F\le G+\norm{F-G}_\infty\)
proves nonexpansiveness.
\end{proof}

The following lemma realizes one coordinate rearrangement up to a prescribed
error.  Its proof, given after the proof of Lemma~\ref{cmp:comparison},
uses smooth families of planar disk maps.

\begin{lemma}[Lifted rearrangement]\label{cmp:lifted}
For every continuous convex \(F:\Delta\to\R\), every coordinate \(j\),
and every \(\eps>0\), there is a Hamiltonian diffeomorphism
\(q_j:V\to V\) such that
\[
 F\circ q_j\le R_jF+\eps.
\]
Its isotopy preserves \(V\), and its Hamiltonians Poisson commute
with every outer-face moment near the corresponding face.
\end{lemma}

\subsection{A fixed point with a uniform gap}

\begin{proof}[Proof of Lemma~\ref{cmp:comparison}]
The continuous convex functions on \(\Delta\) form a complete
metric space in the uniform norm. Since \(-P\) is convex,
Lemma~\ref{cmp:rearrangement-properties} implies that, for
\(0<\lambda<1\), the map
\[
 F\longmapsto R(\lambda F-P)
\]
is a contraction of this space, with constant \(\lambda\).
Let \(K_\lambda\) be its fixed point. Integral preservation gives
\begin{equation}\label{cmp:mean}
 \overline K_\lambda
 =-\frac{\overline P}{1-\lambda}.
\end{equation}
The function \(K_\lambda\) is coordinatewise nondecreasing;
in particular \(K_\lambda(0)=\min_\Delta K_\lambda\).

Choose
\[
 0<d<\min\left\{\overline P,\,
               \inf_{\Delta\setminus\tau\Delta}P\right\}.
\]
We claim that, for \(\lambda\) sufficiently close to one,
\begin{equation}\label{cmp:negative-maximum}
 M_\lambda:=\max_\Delta K_\lambda
 <-\frac d{1-\lambda}.
\end{equation}
Suppose otherwise for such a \(\lambda\).
Then \(\lambda M_\lambda\le M_\lambda+d\), and consequently
\(\lambda K_\lambda-P<M_\lambda\) on the outer region.
For \(p\in\tau\Delta\), write \(p=\tau v\) with \(v\in\Delta\).
Convexity, \(K_\lambda(0)\le\overline K_\lambda\), and
\eqref{cmp:mean} give
\[
 K_\lambda(p)
 \le \tau M_\lambda-
       \frac{(1-\tau)\overline P}{1-\lambda}
 \le M_\lambda-
       \frac{(1-\tau)(\overline P-d)}{1-\lambda}.
\]
It follows that
\[
 \lambda K_\lambda(p)-P(p)
 \le M_\lambda+d
 -\frac{\lambda(1-\tau)(\overline P-d)}{1-\lambda}
 -\min_\Delta P<M_\lambda
\]
when \(\lambda\) is close enough to one. The choices here depend
only on the displayed constants, so the bound is uniform on
both regions. We have shown
\(\max_\Delta(\lambda K_\lambda-P)<M_\lambda\), contradicting
the fixed-point equation and
\(\max RF\le\max F\). This proves
\eqref{cmp:negative-maximum}.

Fix such a \(\lambda\). Since
\((1-\lambda)K_\lambda\le-d\), order preservation and
translation invariance give
\begin{equation}\label{cmp:rearranged-gap}
 R(K_\lambda-P)
 \le R(\lambda K_\lambda-P)-d=K_\lambda-d.
\end{equation}
Set \(F_0=K_\lambda-P\) and \(F_j=R_jF_{j-1}\).
All these functions are continuous and convex.
Use Lemma~\ref{cmp:lifted} to choose \(q_j\) with
\[
 F_{j-1}\circ q_j\le F_j+\eps_j,
 \qquad \sum_j\eps_j<d/2.
\]
In the order \(q=q_1\circ\cdots\circ q_m\), successive
substitution of these inequalities yields
\[
 (K_\lambda-P)\circ q
 \le R(K_\lambda-P)+\sum_j\eps_j
 <K_\lambda-d/2.
\]
Each constituent isotopy preserves \(V\) and preserves every
outer-face moment on a sufficiently small collar.
Choose common threshold collars
\(c_\nu-\epsilon_\nu< b_\nu\cdot p\le c_\nu\)
for the finitely many maps.
Such collars are invariant because the moment is constant
along each flow there. Run the isotopy for \(q_m\) first, then that for \(q_{m-1}\)
on its output, and continue in this order. Each stage uses one of
the original generating Hamiltonians and retains its commutation
with every face moment. Reparametrization makes the concatenated
path smooth and preserves these commutation relations.

Finally approximate \(K_\lambda\) uniformly on \(\Delta\) by
a smooth function \(K(p)\), defined on a neighborhood.
Polynomial approximation is sufficient.
An approximation error smaller than \(d/4\) changes the last
comparison by less than \(d/2\), leaving
\((K-P)\circ q<K\). This proves the lemma.
\end{proof}

\subsection{Realization by Hamiltonian maps}

We now prove the realization lemma used above.  Related constructions
prescribe planar symplectic maps by nested equal-area curves;
see \cite[Lemma~3.2]{SchlenkFolding}.  We need smooth parameter
dependence and relative Hamiltonian control as well.

Write \(D(A)=\{z\in\C:\pi|z|^2<A\}\).
A smooth family of embedded closed disks means a family of smooth
embeddings of the closed unit disk, smooth also in the parameters on
an open neighborhood of their parameter set.

\begin{lemma}[Nested disks]\label{cmp:planar-disks}
Let \(Y\) be a compact convex subset of a Euclidean space, let
\(0<a_1<\cdots<a_l<A\), and let
\[
 E_1(y)\Subset\Int E_2(y)\Subset\cdots
       \Subset\Int E_l(y)\Subset D(A)
 \qquad(y\in Y)
\]
be smooth families of embedded closed disks of areas \(a_1,\ldots,a_l\).
There is a smooth family of Hamiltonian isotopies of \(D(A)\),
supported in a common compact subset, whose endpoint maps carry each
centered disk of area \(a_j\) onto \(E_j(y)\).
\end{lemma}

\begin{proof}
First construct ordinary orientation-preserving diffeomorphisms with
these properties. At a fixed parameter \(y_0\), a disk embedding can
be shrunk within its image, straightened while small, moved through
the ambient disk, and expanded to another prescribed disk. The
velocity along its moving boundary extends to a vector field in a
tubular neighborhood and then, by a cutoff, to a compactly supported
ambient field. The resulting flow extends the disk isotopy.
Apply this construction first to the outermost boundary and then to
each successive boundary inside its already positioned outer disk.
This gives a compactly supported diffeomorphism \(\phi_{y_0}\),
isotopic to the identity, carrying the centered disks to the prescribed
ones.

For other parameters, follow the disk boundaries along the segment
\(y_0+t(y-y_0)\). Their velocities extend in disjoint tubular
neighborhoods of the finitely many boundaries. These extensions can
be chosen smoothly in \((t,y)\), by local extensions and a partition
of unity. Compactness and strict nesting give uniform collars and
support away from the ambient boundary. Integrating the fields
produces diffeomorphisms \(\phi_y\), with ordinary isotopies smooth in
\(y\), carrying all the centered disks to their prescribed images.

Let \(\omega\) be the standard area form and
\(\eta_y=\phi_y^*\omega-\omega\).
This is a compactly supported two-form of integral zero.
Compactly supported primitives can be chosen linearly, and hence
smoothly in \(y\). For completeness, identify the open ambient disk
smoothly with \(\R^2\), write \(\eta=f(x,t)\,\dd x\wedge\dd t\),
and choose a fixed compactly supported function \(\rho\) with
\(\int\rho=1\). Put
\[
 g(t)=\int_\R f(x,t)\,\dd x,\quad
 F(x,t)=\int_{-\infty}^x\bigl(f(v,t)-\rho(v)g(t)\bigr)\,\dd v,
 \quad G(t)=\int_{-\infty}^t g(v)\,\dd v.
\]
Then
\[
 \alpha=F\,\dd t-\rho(x)G(t)\,\dd x,\qquad
 \dd\alpha=\eta.
\]
The zero marginal and zero total integral give compact support, with
a common support for compact families.

For every centered disk \(D_j\) of area \(a_j\),
\[
 \int_{D_j}\eta_y
 =\operatorname{area}E_j(y)-a_j=0.
\]
Thus the restriction of \(\alpha_y\) to \(\partial D_j\) is exact.
Subtract the differentials of functions extending these circle
primitives into disjoint annular collars. The resulting primitive
\(\widehat\alpha_y\) vanishes on the tangent line of each circle.
Apply Moser's method~\cite{Moser1965} to
\(\omega+t\eta_y\), \(0\le t\le1\), using
\(\widehat\alpha_y\). This path consists of positive area forms.
Its vector fields are tangent to the designated circles, since
\(\widehat\alpha_y\) vanishes on their tangents. The resulting
correction \(\chi_y\) preserves every \(D_j\) and satisfies
\(\chi_y^*\phi_y^*\omega=\omega\).
Hence \(\psi_y=\phi_y\circ\chi_y\) is area preserving and has the
required disk images.

It remains to provide a symplectic isotopy with this endpoint.
The construction has already supplied an ordinary compactly
supported isotopy \(\psi_{y,t}\) from the identity to \(\psi_y\).
For each \((y,t)\), apply the same Moser correction, now without
circle conditions, to
\[
 \omega+s(\psi_{y,t}^*\omega-\omega),\qquad 0\le s\le1.
\]
Use the linear primitive operator above. When
\(\psi_{y,t}^*\omega=\omega\), that primitive and the correction are
zero and the identity, respectively. In particular, the corrections
are the identity at both endpoints \(t=0,1\). The corrected isotopy
is therefore symplectic, has endpoint \(\psi_y\), and is smooth in
the parameters. On a disk, a compactly supported symplectic vector
field has a Hamiltonian normalized to vanish on the complementary
component adjacent to the ambient boundary. These Hamiltonians depend
smoothly on the parameters and are supported in a common larger compact
subset of the disk. This yields the required Hamiltonian isotopies.
\end{proof}

\begin{lemma}[Planar rearrangement]\label{cmp:disk-realization}
Let \(f_y:[0,1]\to\R\) be a jointly continuous family of convex
functions, with \(y\) in a compact convex parameter set. Let \(g_y\)
be their increasing rearrangements with respect to Lebesgue measure.
For every \(\eps>0\), there is a smooth family of compactly supported
Hamiltonian isotopies of \(D(1)\), with endpoints \(\psi_y\), such that
\[
 f_y\bigl(\pi|\psi_y(z)|^2\bigr)
 \le g_y(\pi|z|^2)+\eps
 \qquad(z\in\overline{D(1)}).
\]
Their supports lie in a common compact subset of \(D(1)\).
\end{lemma}

\begin{proof}
The minimum formula~\eqref{cmp:interval-formula} shows that
\(g_y(h)\) is jointly continuous in \((y,h)\).

Choose a small \(\eps'>0\) and a sufficiently fine grid
\(0<h_1<\cdots<h_l<1\), including points sufficiently close to both
endpoints. The open intervals
\[
 I_j(y)=
 \{v\in(0,1):f_y(v)<g_y(h_j)+\eps'\}
\]
are nested in \(j\), and each has length greater than \(h_j\).
Indeed, the closed sublevel at \(g_y(h_j)\) contains an interval
of length \(h_j\), and the strict enlargement gives extra room.
Joint continuity and compactness make this room uniform for the
finite grid.

Choose increasing numbers \(w_j<1\), sufficiently close to one, so
that \(h_j/w_j\) is strictly increasing in \(j\) and is smaller than
the available length of \(I_j(y)\) for every \(y\).
There are strictly nested closed intervals \(J_j(y)\), of lengths
\(h_j/w_j\), contained in \(I_j(y)\). To choose them, first choose
the smallest interval and enlarge successively inside the next
available open interval. The admissible centers satisfy strict
affine inequalities expressing containment and nesting. Locally
constant admissible centers remain admissible near a given
parameter. A smooth partition of unity preserves these convex
constraints, and therefore gives centers smooth on a neighborhood
of the parameter set.

In the open polar chart \((v,\theta)\in(0,1)^2\), take the rectangles
\[
 J_j(y)\times
 \left[\frac{1-w_j}{2},\frac{1+w_j}{2}\right].
\]
Their areas are \(h_j\). Round the corners smoothly and adjust the
width in the \(v\)-direction slightly to restore the area exactly.
For example, superellipses of a common sufficiently large even
exponent approximate all these rectangles, and their areas are
corrected by one additional dilation in \(v\).
The strict margins ensure that the resulting smooth closed disks
\(E_j(y)\) remain nested and lie in \(I_j(y)\times(0,1)\).
All choices are uniform because the parameter set and grid are
compact and finite.

Apply Lemma~\ref{cmp:planar-disks} to these disks. If \(v\le h_j\),
the image of a point of radial area \(v\) lies in \(E_j(y)\), and
therefore its \(f_y\)-value is below \(g_y(h_j)+\eps'\).
Use the next larger grid point and uniform continuity of \(g_y\)
to bound this by \(g_y(v)+2\eps'\).
For \(v>h_l\), use
\(f_y\le\max f_y=g_y(1)\) and choose \(h_l\) close enough to one.
Taking \(\eps'\) small proves the assertion. The Hamiltonians
extend by zero across the ambient boundary, so the same inequality
holds on the closed disk.
\end{proof}

\begin{proof}[Proof of Lemma~\ref{cmp:lifted}]
First suppose that a slice has a smooth positive area scale
\(\ell(y)\). Apply Lemma~\ref{cmp:disk-realization} to
\(f_y(v)=F(y,\ell(y)v)\).
If \(B_t(y,z)\) generates the resulting maps on \(D(1)\), then
the Hamiltonian
\[
 \ell(y)B_t\bigl(y,z_j/\sqrt{\ell(y)}\bigr)
\]
on the whole fiber produces the desired scaled disk motion.
The other moments remain fixed because this Hamiltonian has no
dependence on their angles. Their angles may change, which does
not affect a toric function.

To treat the actual endpoint \(L\), choose a small \(\eta>0\) and
work first on the compact convex parameter set
\[
 Y_\eta=\{y:L(y)\ge\eta\}.
\]
Choose a smooth function \(\ell\) strictly below \(L\), uniformly
close to it on \(Y_\eta\), and with \(\ell\ge\eta/2\).
For instance, a smooth soft minimum of the finitely many affine
expressions defining \(L\), shifted downward by a small positive
constant, has these properties when its approximation error is small.
The planar construction and its
Hamiltonians extend smoothly to a neighborhood of \(Y_\eta\).

For \(h\le\ell(y)\), any interval of length \(h\) in \([0,L(y)]\)
can be shifted into \([0,\ell(y)]\) by a distance at most
\(L(y)-\ell(y)\). Hence, for a uniform modulus of continuity
\(\omega_F\),
\begin{equation}\label{cmp:shortened-slice}
 R_{[0,\ell]}F(h)
 \le R_{[0,L]}F(h)+\omega_F(L-\ell).
\end{equation}
For \(\ell<h\le L\), let \([a,a+h]\) be a minimizing interval
for the right-hand rearrangement. Since \(a\le L-h<L-\ell\),
\begin{equation}\label{cmp:outer-strip}
 F(y,h)\le R_{[0,L]}F(h)+\omega_F(L-\ell).
\end{equation}
Thus the identity map in this outer strip has the required
one-sided estimate. The planar Hamiltonians vanish near their
rescaled boundary, so this identity extension is smooth.

Choose a smooth cutoff \(0\le\chi(y)\le1\), supported where
\(L>\eta\) and equal to one where \(L\ge2\eta\), with its support
contained in the neighborhood where the preceding choices are
defined. The full lifted Hamiltonian is
\begin{equation}\label{cmp:cutoff-hamiltonian}
 \chi(y)\ell(y)
 B_{\chi(y)t}\bigl(y,z_j/\sqrt{\ell(y)}\bigr),
 \qquad 0\le t\le1,
\end{equation}
extended by zero elsewhere. At fixed \(y\), it runs the scaled
disk isotopy to time \(\chi(y)\). On slices of length at most
\(2\eta\), any slice-preserving motion incurs error at most
\(\omega_F(2\eta)\), since all slice values of \(F\) have that
oscillation bound. Elsewhere \(\chi=1\), and
\eqref{cmp:shortened-slice}--\eqref{cmp:outer-strip}, together
with the planar estimate, apply. Choose \(\eta\), the scale
error, and the planar error in this order to make the total
error less than \(\eps\).

All lifts are smooth at the other coordinate axes because
the other moments are smooth functions of the complex variables.
The cutoff removes the possible degeneracy at \(L=0\).
To check the outer faces, suppose first that \(b_{\nu j}>0\).
Equality on that face forces \(p_j=L(y)\). On active slices the
common compact planar support satisfies
\[
 p_j\le(1-\delta)\ell(y)
\]
for some \(\delta>0\), while \(\ell(y)\ge\eta/2\).
Writing \(L_\nu(y)\) for this face's affine endpoint, its support
therefore satisfies
\[
 c_\nu-b_\nu\cdot p
 =b_{\nu j}(L_\nu(y)-p_j)
 \ge b_{\nu j}\delta\ell(y)
 \ge b_{\nu j}\delta\eta/2>0.
\]
Thus the Hamiltonian vanishes on a genuine uniform neighborhood
of that face. If \(b_{\nu j}=0\), the Hamiltonian commutes with
\(b_\nu\cdot p\) everywhere. These statements imply that the
flow preserves all defining inequalities of \(V\) and gives
the stated face property.
\end{proof}
